\begin{frame}{Раздел 4: Дизайн системы и оптимизация}
    На основе анализа проблем (сбои, задержки оценки) разработаны две системы:
    \begin{enumerate}
        \item \textbf{Fault-Tolerant Pretraining (Отказоустойчивый претренинг)}
        \item \textbf{Decoupled Scheduling for Evaluation (Декомпозированное планирование для оценки)}
    \end{enumerate}
    Обе интегрированы в фреймворк \texttt{InternEvo}.
\end{frame}

\begin{frame}{4.1 Fault-Tolerant Pretraining: Компоненты}
    \begin{itemize}
        \item \textbf{Асинхронное чекпойнтинг:}
            \begin{itemize}
                \item Ускоряет сохранение состояния модели в \textbf{3.6--58.7x} раз по сравнению с синхронным подходом.
                \item Позволяет делать чекпойнты часто, не блокируя обучение.
            \end{itemize}
        \item \textbf{Диагностика с помощью LLM:}
            \begin{itemize}
                \item Комбинация эвристических правил и LLM для анализа логов.
                \item LLM помогает классифицировать и находить первопричины редких или сложных ошибок.
            \end{itemize}
        \item \textbf{Холлистический детектор:}
            \begin{itemize}
                \item Набор инструментов для точного определения неисправного узла (GPU, сеть, память).
                \item Автоматический перезапуск задачи с последнего корректного чекпойнта.
            \end{itemize}
        \item \textbf{Результат:} Снижение необходимости ручного вмешательства и простоев.
    \end{itemize}
\end{frame}

\begin{frame}{4.2 Decoupled Scheduling for Evaluation}
    \begin{itemize}
        \item \textbf{Проблема:} Длительные простои GPU из-за загрузки модели с удаленного хранилища и CPU-вычислений метрик.
    \end{itemize}
    \textbf{Решения:}
    \begin{itemize}
        \item \textbf{Декомпозиция загрузки модели:} Кэширование моделей локально на узлах, чтобы избежать повторной загрузки по сети для каждой задачи.
        \item \textbf{Декомпозиция вычисления метрик:} Вынос тестов синтезированных программ (CPU-bound) на отдельные пулы ресурсов, освобождая GPU для инференса.
        \item \textbf{Балансировка на основе знаний о датасетах:} Использование prior knowledge (размер датасета, сложность) для равномерного распределения подзадач по GPU.
    \end{itemize}
    \textbf{Результат:} Сокращение общего времени выполнения оценки (\textit{ makespan}) до \textbf{1.8x}.
\end{frame}